\preface
\thispagestyle{empty}

Mathematics begins with intuition, but it becomes reliable only when its assumptions and arguments are made explicit. This manuscript grew from an attempt to understand that transition: how familiar calculations are reorganized into definitions, theorems, proofs, algorithms, and models that can be reused in unfamiliar settings.

The book is best described as a \emph{guided second-pass reference}. It is written primarily for university students who have already encountered at least some calculus and linear algebra and now want to connect those courses with abstract algebra, topology, complex analysis, graph theory, and probability. It is not intended to replace a first textbook in every subject. Instead, it supplies a coherent map, substantial proofs of central results, worked examples, applications, and exercises that help a reader move between computational skill and structural understanding.

The opening chapters on logic and set theory are optional foundations, not prerequisites for every reader. A student whose immediate goal is calculus may begin with Mathematical Analysis I. A reader interested in computation and data may begin with Linear Algebra, then move to Graph Theory and Probability. The foundational chapters remain available for readers who want to examine how mathematical language and constructions are formalized, but the architecture of the book no longer requires every reader to pass through ZFC before doing analysis.

The level of detail varies because the subjects serve different roles. Analysis and linear algebra receive extended development because they provide techniques used throughout the rest of the book. Later chapters are selective introductions, but they are not intended to be theorem catalogues. Each chapter now states learning objectives, develops representative worked examples, connects the theory to scientific or computational practice, and ends with exercises and selected guidance.

Proofs are included whenever the mechanism of the argument is central to the subject. In particular, limit arguments, compactness arguments, quotient constructions, convergence theorems, and probabilistic truncation methods are written out rather than hidden behind the phrase “standard proof.” Results whose complete proofs require machinery far beyond the scope of the chapter are explicitly identified as such. This distinction is important: a concise proof is not the same as a missing proof, and an advanced theorem should not be presented as elementary merely because its statement is short.

Applications are used to clarify structure rather than to decorate it. Linear algebra is connected to least squares, data compression, and principal component analysis; abstract algebra to coding and cryptography; multivariable analysis to optimization and physical fields; topology to continuity of data representations and quotient constructions; complex analysis to signal processing and potential theory; graph theory to routing, networks, and graph-based learning; and probability to statistical inference, simulation, and stochastic processes. These examples are deliberately modest: their purpose is to show how the mathematical language enters a real model, not to replace a specialized applied course.

The exposition distinguishes personal interpretation from mathematical convention. Phrases such as “the axiom of linear algebra” are replaced by more precise claims: the vector-space axioms are one especially powerful organizing viewpoint, but computational, geometric, and operator-theoretic approaches are equally legitimate. Readers are encouraged to compare formulations and to ask which viewpoint is most effective for a given problem.

A productive way to use this book is to read with pencil and paper. Pause after a definition and construct examples and non-examples. Before reading a proof, identify the hypotheses and predict where each one will be used. Work the chapter exercises without looking immediately at the hints. For computational topics, verify a small case numerically. Mathematics becomes durable when an argument can be reconstructed rather than merely recognized.

This remains a developing manuscript. Its purpose is not to present one final route through undergraduate mathematics, but to provide a precise and usable route whose choices are visible to the reader.

\begin{flushright}
Jinshuo Li \\
Shanghai Jiao Tong University \\
2026.7.3, in Shanghai
\end{flushright}
